\documentclass[11pt,a4paper]{article}

% ===== Paquetes básicos =====
\usepackage[spanish]{babel}
\usepackage[utf8]{inputenc}
\usepackage[T1]{fontenc}
\usepackage{lmodern}
\usepackage{url}
\usepackage{hyperref}
\usepackage{geometry}
\usepackage{setspace}
\usepackage{listings}
\usepackage{xcolor}

\geometry{margin=2.5cm}
\setstretch{1.1}

% Configuración para código
%\lstset{
%	basicstyle=\ttfamily\small,
%	breaklines=true,
%	showstringspaces=false,
%	columns=fullflexible
%}


\lstdefinelanguage{JavaScript}{
	% Palabras clave básicas
	morekeywords=[1]{
		break, case, catch, class, const, continue, debugger, default, delete, do,
		else, export, extends, finally, for, function, if, import, in, instanceof,
		let, new, return, super, switch, this, throw, try, typeof, var, void, while,
		with, yield
	},
	% Literales y tipos
	morekeywords=[2]{
		false, null, true, undefined,
		boolean, number, string, symbol, bigint
	},
	% Objetos y funciones globales
	morekeywords=[3]{
		Array, Boolean, Date, Error, EvalError, Function, JSON, Map, Math,
		Number, Object, Promise, Proxy, RangeError, ReferenceError, RegExp, Set,
		String, SyntaxError, TypeError, URIError, WeakMap, WeakSet,
		console, window, document
	},
	sensitive=true,                % distingue mayúsculas
	% Comentarios
	morecomment=[l]{//},
	morecomment=[s]{/*}{*/},
	morecomment=[s]{/**}{*/},      % estilo JSDoc
	% Cadenas
	morestring=[b]',
	morestring=[b]",
	morestring=[b]`                % template strings
}[keywords, comments, strings]



\lstdefinestyle{jsstyle}{
	language=JavaScript,
	basicstyle=\ttfamily\small,
	keywordstyle=[1]\color{blue}\bfseries,
	keywordstyle=[2]\color{orange!80!black},
	keywordstyle=[3]\color{teal!70!black},
	stringstyle=\color{green!50!black},
	commentstyle=\color{gray},
	numberstyle=\tiny\color{gray},
	numbers=left,
	stepnumber=1,
	breaklines=true,
	showstringspaces=false,
	columns=fullflexible
}


\lstdefinestyle{terminalstyle}{
	basicstyle=\ttfamily\small,
	backgroundcolor=\color{black!5},
	keywordstyle=\color{black},
	commentstyle=\color{black},
	stringstyle=\color{black},
	showstringspaces=false,
	breaklines=true,
}



\title{Monolito m\'{\i}nimo con Node.js y Express}
\author{Gu\'{\i}a de laboratorio}
\date{}

\begin{document}
	
	\maketitle
	
	A continuación se presenta el ``monolito m\'{\i}nimo'' paso a paso, indicando las versiones m\'{\i}nimas necesarias de Node.js y npm.
	%, y el c\'odigo completo en un solo archivo \texttt{index.js}. \cite{web48,web49,web63,web65,web67}
	
	\section{Versiones m\'{\i}nimas requeridas}
	
	Para evitar errores de dependencias o funcionamiento incorrecto, se requieren las siguientes versiones:
	
	\begin{itemize}
		\item \textbf{Node.js (m\'{\i}nimo}: \texttt{16.9.0}. \textbf{recomendado:} LTS, por ejemplo \texttt{18.x} o \texttt{20.x}. )
		\item \textbf{npm}: la versi\'on que venga incluida con la instalaci\'on de Node LTS elegida (para Node 18/20 normalmente es 8.x--10.x).
	\end{itemize}
	
	Para comprobar las versiones en la terminal:
	
	\begin{lstlisting}[style=terminalstyle]
		node -v   # deberia ser >= 16.9.0 (ideal: 18.x o 20.x)
		npm -v
	\end{lstlisting}
	
	\section{Crear el proyecto}
	
	\begin{enumerate}
		\item Crear la carpeta del proyecto:
		
		\begin{lstlisting}[style=terminalstyle]
			mkdir tienda-monolito
			cd tienda-monolito
		\end{lstlisting}
		
		\item Inicializar el proyecto Node:
		
		\begin{lstlisting}[style=terminalstyle]
			npm init -y
		\end{lstlisting}
		
		\item Instalar Express:
		
		\begin{lstlisting}[style=terminalstyle]
			npm install express
		\end{lstlisting}
		
		Esto instalar\'a una versi\'on reciente de Express compatible con una versi\'on de Node. 
	\end{enumerate}
	
	\section{Crear \texttt{index.js} con el monolito m\'{\i}nimo}
	
	En la carpeta \texttt{tienda-monolito}, crear el archivo \texttt{index.js} con el siguiente contenido completo:
	
	\begin{lstlisting}[style=jsstyle]
		// index.js
		// Monolito minimo con Express
		
		const express = require('express');
		
		const app = express();
		const PORT = 3000;
		
		// Middleware para poder leer JSON en el body de las peticiones
		app.use(express.json());
		
		// ===== Datos en memoria =====
		
		// Catalogo de productos (normalmente esto estaria en una BD)
		const productos = [
		{ id: 1, nombre: 'Laptop basica', precio: 12000 },
		{ id: 2, nombre: 'Mouse inalambrico', precio: 350 },
		{ id: 3, nombre: 'Teclado mecanico', precio: 1500 }
		];
		
		// Pedidos registrados durante la ejecucion
		const pedidos = [];
		
		// ===== Rutas =====
		
		// GET /productos -> devuelve el listado de productos
		app.get('/productos', (req, res) => {
			res.json(productos);
		});
		
		// POST /pedidos -> registra un nuevo pedido en memoria
		app.post('/pedidos', (req, res) => {
			const nuevoPedido = req.body;
			
			// Validacion muy basica
			if (
			!nuevoPedido ||
			!nuevoPedido.productos ||
			!Array.isArray(nuevoPedido.productos)
			) {
				return res.status(400).json({
					mensaje: 'El pedido debe incluir un arreglo de productos.'
				});
			}
			
			// Asignar un id incremental
			const id = pedidos.length + 1;
			const pedidoConId = { id, ...nuevoPedido };
			
			// Guardar en memoria
			pedidos.push(pedidoConId);
			
			// Responder con el pedido creado
			res.status(201).json(pedidoConId);
		});
		
		// (Opcional) GET /pedidos -> ver todos los pedidos registrados
		app.get('/pedidos', (req, res) => {
			res.json(pedidos);
		});
		
		// ===== Arrancar el servidor =====
		
		app.listen(PORT, () => {
			console.log(`Servidor escuchando en http://localhost:${PORT}`);
		});
	\end{lstlisting}
	
	
	Este archivo \'unico contiene datos en memoria, rutas HTTP y el arranque del servidor, por lo que constituye un ``monolito m\'{\i}nimo'' funcional.
	
	\section{Ejecutar el servidor}
	
	Desde la carpeta del proyecto, ejecutar:
	
	\begin{lstlisting}[style=terminalstyle]
		node index.js
	\end{lstlisting}
	
	Si todo est\'a bien, aparecer\'a en la terminal:
	
	\begin{lstlisting}[style=terminalstyle]
		Servidor escuchando en http://localhost:3000
	\end{lstlisting}
	
	El proceso se mantendr\'a en ejecuci\'on, indicando que el servidor est\'a levantado. 
	
	\section{Probar las rutas}
	
	\subsection{GET /productos}
	
	Probar en el navegador o Postman:
	
	\begin{lstlisting}[style=terminalstyle]
		http://localhost:3000/productos
	\end{lstlisting}
	
	Se deber\'ia mostrar el arreglo de productos en formato JSON. 
	
	\subsection{POST /pedidos}
	
	Ejemplo usando \texttt{curl} (tambi\'en puede hacerse con Postman):
	
	\begin{lstlisting}[style=terminalstyle, caption={Usando linux, Mac}]
		curl -X POST http://localhost:3000/pedidos \
		-H "Content-Type: application/json" \
		-d "{
			\"cliente\": \"Ana Perez\",
			\"productos\": [1, 3],
			\"total\": 13500
		}"
	\end{lstlisting}
	
	
	\begin{lstlisting}[style=terminalstyle, caption={Usando cmd de Windows}]
		curl -X POST "http://localhost:3000/pedidos" ^
		-H "Content-Type: application/json" ^
		-d "{
			\"cliente\": \"Ana Perez\",
			\"productos\": [1, 3],
			\"total\": 13500
		}"
	\end{lstlisting}

	\begin{lstlisting}[style=terminalstyle, caption={Usando PowerShell de Windows}]
	curl -X POST "http://localhost:3000/pedidos" `
	-H "Content-Type: application/json" `
	-d '{
		"cliente": "Ana Perez",
		"productos": [1, 3],
		"total": 13500
	}'
\end{lstlisting}	
	
	
	La respuesta esperada es un JSON similar a:
	
	\begin{lstlisting}[style=terminalstyle]
		{
			"id": 1,
			"cliente": "Ana Perez",
			"productos": [1, 3],
			"total": 13500
		}
	\end{lstlisting}
	
	\subsection{GET /pedidos (opcional)}
	
	Para ver todos los pedidos registrados durante la ejecuci\'on:
	
	\begin{lstlisting}[style=terminalstyle]
		http://localhost:3000/pedidos
	\end{lstlisting}
	
	Esto mostrar\'a el arreglo \texttt{pedidos} con todos los pedidos creados en memoria.
	
%	\section*{6. Comentario final}
%	
%	Con Node \(\geq 16.9.0\) (idealmente una versi\'on LTS como 18 o 20) y este archivo \texttt{index.js}, la pr\'actica funciona sin el error de \texttt{Object.hasOwn} y sirve como base para discutir monolito, API REST y futuras refactorizaciones hacia servicios m\'as peque\~nos. \cite{web48,web49,web63,web65,web67}
%	
%	\bibliographystyle{plain}
%	\begin{thebibliography}{9}
%		
%		\bibitem{web48}
%		GitHub.
%		\newblock \emph{TypeError: Object.hasOwn is not a function · Issue \#41471}.
%		\newblock 2022.
%		
%		\bibitem{web49}
%		Stack Overflow.
%		\newblock \emph{TypeError: Object.hasOwn is not a function - javascript}.
%		\newblock 2023.
%		
%		\bibitem{web63}
%		MDN Web Docs.
%		\newblock \emph{Object.hasOwn()}.
%		\newblock 2025.
%		
%		\bibitem{web65}
%		GitHub.
%		\newblock \emph{[Discussion] Upgrade node and npm to latest LTS versions}.
%		\newblock 2023.
%		
%		\bibitem{web71}
%		npm.
%		\newblock \emph{node-lts-versions}.
%		\newblock 2025.
%		
%		\bibitem{web67}
%		npm.
%		\newblock \emph{express} (version info).
%		\newblock 2024.
%		
%		\bibitem{web73}
%		npm.
%		\newblock \emph{express} (versions overview).
%		\newblock 2025.
%		
%		\bibitem{web18}
%		Express.js.
%		\newblock \emph{Hello world example}.
%		\newblock 2024.
%		
%		\bibitem{web19}
%		Express.js.
%		\newblock \emph{Express "Hello World" example}.
%		\newblock 2024.
%		
%		\bibitem{web23}
%		JSON Parser Blog.
%		\newblock \emph{How to Parse JSON Request Bodies in Express.js with Body-Parser}.
%		\newblock 2025.
%		
%		\bibitem{web22}
%		JavaScript Tutorial.
%		\newblock \emph{Building the Express Hello World App}.
%		\newblock 2024.
%		
%		\bibitem{web25}
%		Tutorialspoint.
%		\newblock \emph{ExpressJS - Hello World!}.
%		\newblock n.d.
%		
%	\end{thebibliography}
	
\end{document}
